%META {"question":"Tam giác ABC (nhọn). Đoạn thẳng BH có là đường cao của tam giác ABC không?","options":["Có — vì BH ⊥ AC và đi từ đỉnh B xuống cạnh đối diện AC","Không — vì H không là trung điểm của AC","Không — vì BH không đi qua đỉnh A","Không — vì BH không nằm trong tam giác"],"answer":0,"note":"Đường cao từ đỉnh B phải vuông góc với cạnh đối diện AC. H nằm trên AC và BH ⊥ AC (dấu góc vuông tại H) → BH là đường cao ứng với cạnh AC."}
\documentclass[border=5pt]{standalone}
\usepackage{tkz-euclide}
% Màu dùng chung cho toàn bộ hình vẽ (ảnh stack với nền tối của web)
\definecolor{tminc}{RGB}{226,232,240}   % chữ + nét chính
\definecolor{tmblue}{RGB}{0,210,255}    % nhấn neon xanh
\definecolor{tmpink}{RGB}{255,0,200}    % nhấn neon hồng
\definecolor{tmgreen}{RGB}{0,255,136}   % nhấn neon xanh lá
\color{tminc}

\begin{document}
\begin{tikzpicture}[line width=1.1pt]
  \coordinate (A) at (0,0);
  \coordinate (B) at (3.8,0);
  \coordinate (C) at (2.2,2.8);
  \tkzDefPointBy[projection=onto A--C](B) \tkzGetPoint{H}
  \tkzDrawPolygon[tmblue,line width=1.1pt](A,B,C)
  \draw[tmpink] (B) -- (H);
  \tkzMarkRightAngles[size=0.28,color=tmpink,line width=1pt](B,H,A)
  \tkzDrawPoints[size=4.5,color=tmpoint,fill=tmpoint](A,B,C,H)
  \tkzLabelPoint[below left=1pt](A){$A$}
  \tkzLabelPoint[below right=1pt](B){$B$}
  \tkzLabelPoint[above right=1pt](C){$C$}
  \tkzLabelPoint[above=3pt](H){$H$}
\end{tikzpicture}
\end{document}
